\section{Categorization of Memory-Use Failures}
\label{sec:problem}

Persistent memory affects an agent through the information that is selected from memory and used to handle a new request. We use
\(
U(q,M)
\)
to denote the observable memory-use behavior of an agent when processing query $q$ with memory state $M$. Depending on the system, $U(q,M)$ may
include the memories retrieved for the query, their ranking, or the
memory-dependent response produced by the agent. We do not assume that a
memory-use failure is caused by any particular internal component.

We categorize memory-use failures by changing one of the two inputs to
$U(q,M)$ while keeping the other fixed. This gives two categories:

\textbf{I. Query-related failures} occur when the memory state remains fixed,
but memory use does not correctly follow an adaptation from $q$ to
$q'$. The expected behavior depends on the relation between the two
queries. For example, if $q'$ is a paraphrase of $q$, the same relevant
memory should continue to support the request; if $q'$ asks for unsupported
information, the agent should not rely on unrelated memory. We write this
requirement as
\fbox{\(
\mathcal{C}_Q\!\left(U(q,M), U(q',M); q,q'\right)
\)},
where $\mathcal{C}_Q$ specifies the relation expected from the query
change, such as preserving memory use under equivalent queries or avoiding
unsupported memory use for an unanswerable query. A query-related failure
occurs when this condition is violated.

\textbf{II. Memory-state failures} occur when the query remains fixed, but
memory use does not correctly follow a controlled change from $M$ to
$M'$. A relevant update or deletion should be reflected when the agent later
uses memory, while an unrelated change should not disturb
the information used for the same request. We write this requirement as
\fbox{\(
\mathcal{C}_M\!\left(U(q,M), U(q,M'); M,M'\right)
\)},
where $\mathcal{C}_M$ specifies the relation expected from the memory
change, such as preferring the updated value after a replacement or no
longer using an entry after it has been deleted. A memory-state failure
occurs when this condition is violated.

This categorization is based on the change that triggers the failure, thus agnostic to how the memory system is integrated to LLM agents. It therefore applies
to systems with explicit retrieval interfaces as well as systems whose
memory use mechanism is hidden. The mutation operators in
Section~\ref{sec:method} instantiate these two categories with concrete
query and memory changes.
